\section{Round-nine reconstruction: a physical impact section and a directional current calculus}
\label{sec:r9-a1}

The physical Poincar\'e section and the symbolic natural extension are different
objects.  The former is a two-dimensional symplectic section with a finite
one-step singularity set; the latter is a Cantor path space used to resolve all
iterated one-sided limits.  No symbolic graph is called a Liouville section in
this revision.

\subsection{The regular physical section and its path factor}

Let $p_i(a)>0$, $\sum_{i=1}^d p_i(a)=1$, and
$s_i(a)=\sum_{j\le i}p_j(a)$.  On
$\Sigma=(0,1)^2$ put
\[
 I_i(a)=(s_{i-1}(a),s_i(a)),
 \qquad
 B_a(q,p)=\left(
 \frac{q-s_{i-1}(a)}{p_i(a)},
 s_{i-1}(a)+p_i(a)p
 \right),\quad q\in I_i(a).
\]
The one-step seam set is $\mathscr S_a=\bigcup_i\partial I_i(a)\times(0,1)$.
Let
\[
 X_a=\Sigma\setminus\bigcup_{n\in\mathbb Z}B_a^{-n}\mathscr S_a.
\]
It has full Lebesgue measure and is invariant.  Put
$\Omega=\{1,\ldots,d\}^{\mathbb Z}$ and let $\nu_a$ be the Bernoulli law with
weights $p(a)$.

\begin{theorem}[Physical--symbolic separation and exact factor]
\label{thm:r9-a1-factor}
The restriction $B_a:X_a\to X_a$ is an invertible measure-preserving map.
The itinerary map
\[
 \mathfrak c_a:X_a\longrightarrow\Omega,
 \qquad
 (\mathfrak c_a x)_n=i\ \Longleftrightarrow\ B_a^nx\in I_i(a)\times(0,1),
\]
is a bimeasurable isomorphism modulo null sets and
$(\mathfrak c_a)_\#m=\nu_a$.  Endpoint expansions belong only to the symbolic
completion; they are not extra points of the physical symplectic section.
\end{theorem}

\begin{proof}
Each branch has Jacobian one and sends a vertical strip bijectively to the
corresponding horizontal strip.  Removing all forward and backward seam
preimages makes every iterate and inverse iterate single valued.  A cylinder
$[i_{-r}\ldots i_s]$ pulls back to a rectangle of area
$\prod_{k=-r}^sp_{i_k}(a)$.  Cylinders generate the two sigma fields, and their
diameters tend to zero in both stable and unstable directions.  Hence the
itinerary is one-to-one and onto modulo endpoint sequences and carries
Lebesgue measure to $\nu_a$.
\end{proof}

\subsection{An exact Hamiltonian impact suspension}

For branch $i$ let
\[
 S_i^a(q,P)=\frac{(q-s_{i-1})(P-s_{i-1})}{p_i}.
\]
Then $Q=\partial_PS_i^a$ and $p=\partial_qS_i^a$, so the branch reset is exact
symplectic.  Consider the extended phase space
\[
 \mathcal M_a=\Sigma\times\mathbb T_\tau\times\mathbb R_E
\]
with symplectic form $dq\wedge dp+d\tau\wedge dE$.  Between impacts use
$H=E$.  At $\tau=1$ apply the canonical reset generated by $S_i^a$ when the
incoming point lies in $I_i(a)\times(0,1)$, and reset $\tau$ to zero.  The
incoming side determines the reset at a wall; the wall itself has zero
Liouville flux.

\begin{theorem}[Autonomous Hamiltonian impact realization]
\label{thm:r9-a1-impact}
The preceding data define a piecewise smooth autonomous Hamiltonian impact
flow whose first-return time to $\{\tau=0\}$ is one and whose return map is
$B_a$ on the full-measure invariant set $X_a$.  The reset preserves the
section form $dq\wedge dp$ and the induced Liouville flux.  The construction is
$C^{m+1}$ in $a$ on every compact parameter interval on which
$\min_i p_i(a)>0$.
\end{theorem}

\begin{proof}
Hamilton's equations give $\dot\tau=1$, $\dot q=\dot p=0$ between impacts.
At the impact the generating-function identities give exactly the displayed
branch formula.  The determinant is one and the canonical one-form changes by
$dS_i^a$, hence the reset is exact symplectic.  Branch interiors are disjoint,
so the impact rule is deterministic.  Seam points have zero incoming flux and
are removed in $X_a$.  Smooth parameter dependence follows from the explicit
formula and the uniform lower bound on the branch widths.
\end{proof}

This is a Hamiltonian impact system, not a claim that the Cantor natural
extension is a smooth section.  A smooth regularization is obtained by
replacing the instantaneous reset by a thin canonical layer generated by
$S_i^a$ on each incoming port; its scattering map converges in every fixed
$C^r$ norm away from the seam.

\subsection{A directional two-sided anisotropic scale}

Let $\mathcal F_{r,s}$ be the sigma field generated by
$\omega_{-r},\ldots,\omega_s$.  Write $\Delta_{r,s}$ for the double martingale
difference
\[
 \Delta_{r,s}f=E(f\mid\mathcal F_{r,s})
 -E(f\mid\mathcal F_{r-1,s})
 -E(f\mid\mathcal F_{r,s-1})
 +E(f\mid\mathcal F_{r-1,s-1}),
\]
with the obvious one-sided conventions.  Fix
$0<\beta<\alpha<1$ and $0<\beta_-<\alpha_-<1$ with
$\beta/\alpha<1$.  Define
\[
 \|f\|_+=|\nu_a(f)|+
 \sum_{r,s\ge0}\alpha^{-r}\beta^{-s}
 \|\Delta_{r,s}f\|_{L^1(\nu_a)},
\]
\[
 \|f\|_-=|\nu_a(f)|+
 \sum_{r,s\ge0}\alpha_-^{-r}\beta_-^{-s}
 \|\Delta_{r,s}f\|_{L^1(\nu_a)}.
\]
Choose the weak parameters so that
$\alpha_-^{-1}<\alpha^{-1}$ and $\beta_-^{-1}<\beta^{-1}$.
Let $\mathscr B_a^+\hookrightarrow\mathscr B_a^-$ be the corresponding
completions.  The transfer operator is
$\mathcal L_af=f\circ\sigma^{-1}$.

\begin{lemma}[Directional Lasota--Yorke estimate]
\label{lem:r9-a1-ly}
There are $C<\infty$ and $\rho=\beta/\alpha<1$ such that
\[
 \|\mathcal L_a^nf\|_+
 \le C\rho^n\|f\|_+ + C\|f\|_-,
 \qquad n\ge0.
\]
The inclusion $\mathscr B_a^+\hookrightarrow\mathscr B_a^-$ is compact.
\end{lemma}

\begin{proof}
Composition with $\sigma^{-1}$ sends a block with past/future depths $(r,s)$
to $(r+1,s-1)$ while $s\ge1$.  Its strong weight is therefore multiplied by
$\beta/\alpha$.  After $s$ steps the block has no future coordinate; those
terms are bounded by the weak norm.  Splitting the martingale sum into
$s\ge n$ and $s<n$ gives the estimate.  The translated-cylinder counterexample
is absorbed because a cylinder at future depth $N$ has initial strong weight
$\beta^{-N}$.  For compactness, truncate $(r,s)$ to a finite rectangle; the
truncation is finite dimensional and the remaining tail is bounded by the
strict ratios $\alpha/\alpha_-$ and $\beta/\beta_-$.
\end{proof}

\begin{theorem}[Spectrum on the directional scale]
\label{thm:r9-a1-spectrum}
The eigenvalue one is simple on $\mathscr B_a^+$ and the remainder of the
spectrum lies in a disk of radius strictly smaller than one.  The reduced
resolvent
\[
 \mathcal R_a=(I-\mathcal L_a)^{-1}(I-\Pi_a)
\]
is bounded from $\mathscr B_a^+$ to itself and from each order of the
parameter-current scale below to the next weaker order.
\end{theorem}

\begin{proof}
Lemma~\ref{lem:r9-a1-ly} gives quasicompactness.  If
$\mathcal L_af=e^{it}f$ with $|e^{it}|=1$, equality in the martingale
contraction forces every future martingale difference to vanish.  Applying the
same argument to the inverse on the dual scale removes every past difference.
Thus $f$ is constant and $e^{it}=1$.  The resolvent statement follows from the
spectral decomposition and the same strong/weak estimate.
\end{proof}

\subsection{All-depth physical currents and response}

For a finite future word $w=i_0\ldots i_n$, the corresponding physical cut is
an affine vertical line whose coefficient is a product of branch widths.  Let
$[\partial w]$ be its oriented line current.  For $k\le m$ define
\[
 \|J\|_{m,\mathrm{cur}}
 =\sum_{k\le m}\sum_{w}
 (1+|w|)^m\,\nu_a([w])^\gamma
 |J_k(w)|,
 \qquad 0<\gamma<1.
\]
On a compact parameter interval,
$\nu_a([w])\le(1-c)^{|w|}$, so the current series and all finite parameter
orders are summable.  The coding map is used only to index these physical
line currents.

\begin{lemma}[Material derivative on every generated cut]
\label{lem:r9-a1-current}
For each finite order $m$, differentiation of the physical transfer operator
has the form
\[
 D_a\mathcal L_a=\bigl(D_a\mathcal L_a\bigr)_{\rm bulk}
 +\sum_w c_w(a)[\partial w],
\]
where the coefficient and all derivatives through order $m-1$ satisfy the
current norm above.  The same statement holds for iterated derivatives, with
corners represented by products of compatible one-dimensional currents.
\end{lemma}

\begin{proof}
Differentiate the branch formula on each finite cylinder before taking the
limit.  The derivative of an indicator is the oriented boundary current and
the derivative of its position is a sum of at most $|w|$ logarithmic branch
width derivatives.  Thus the order-$k$ coefficient is bounded by
$C_k(1+|w|)^k\nu_a([w])$.  Exponential cylinder decay dominates every
polynomial and proves normal convergence.  Compatibility under refinement is
exactly finite additivity of oriented currents, so the completion is
Hausdorff.
\end{proof}

\begin{theorem}[Finite-order physical response]
\label{thm:r9-a1-response}
Let $G_a$ be a physical observable whose pullback to $\Omega$ and whose shape
currents have $m$ finite orders in the preceding scales.  Then
$a\mapsto\int_{X_a}G_a\,dm$ is $C^m$.  Every derivative is a finite sum of
normally convergent words in
\[
 D_a^j\mathcal L_a,
 \qquad \mathcal R_a,
 \qquad D_a^jG_a,
\]
and every seam contribution is the pushforward of an explicitly indexed
physical current.  Common-root finite differences converge to the same
response array.
\end{theorem}

\begin{proof}
Differentiate the eigenrelation
$\nu_a\mathcal L_a=\nu_a$ on the directional scale and solve the centered
part with $\mathcal R_a$.  Lemma~\ref{lem:r9-a1-current} supplies the shape
derivatives and their normal summability.  Iteration gives the usual finite
Kato words with one strong-order loss at each current insertion.  Under the
common inverse-CDF root coupling, the only discontinuities in the parameter
are the same moving cylinder boundaries; dominated convergence away from them
and the current formula on them identify the finite-difference limit with the
operator derivative.
\end{proof}

\subsection{Work and conditional calibration}

Put a one-form $\eta=c_i\,d\tau$ on the complete incoming port of branch $i$.
The work in $n$ returns is $A_n=\sum_{j<n}c_{\omega_j}$.  The canonical tilt is
therefore the exact path-law change
\[
 \frac{d\nu_{a,\theta}}{d\nu_a}\Big|_{\mathcal F_n}
 =\exp\{\theta A_n-n\Lambda_a(\theta)\}.
\]

\begin{theorem}[Mechanical work and conditional coefficient identification]
\label{thm:r9-a1-calibration}
The impact suspension accumulates exactly $A_n$ and its tilted path law is the
Bernoulli exponential family above.  A cash-additive, time-consistent
certainty equivalent has the same coefficient as the mechanical conjugate
field precisely when its likelihood cocycle is required to equal this
canonical mechanical cocycle.  No identification is asserted for an
unrelated preference family.
\end{theorem}

\begin{proof}
Each flight crosses exactly one incoming port and the integral of $c_i d\tau$
over that unit flight is $c_i$.  The cylinder Radon--Nikodym derivative is the
displayed exponential.  Exponential-utility rigidity gives a constant
certainty-equivalent coefficient; equality of the two nonconstant likelihood
cocycles forces equality of their exponents.
\end{proof}
